设 \(\widehat r_{\rm diff}\ge0\) 为具有偶奇偶性且 \(a_{1p}\ne a_{2p}\) 的素数的
实际对数范数。表中给出
\(\widehat R=\widehat s+\widehat r_{\rm diff}\)，因此
\begin{equation}\label{eq:fixed-cube-count}
 \widehat c=\widehat\ell_{\rm pair}+\widehat R/2-\widehat j
                  -\widehat r_{\rm diff}/2
 \le\ell+R/2-j+\tfrac52\eta.
\end{equation}
以此实际上界进行的理想计数在保留 \(J\) 的同时计数固定的 \(q_0\)。
对于一个同时也计数 \(J\) 的对角界，
其完整窗口的指数至多为 \(j+\eta\)，给出合并的
上界计数 \(Z^{\ell+R/2+7\eta/2+\epsilon}\)。\(\mathfrak A_i\) 的
选取以及 \(d_k\mid C\mathfrak B/R_1\) 的选取在 \(b_1,b_2,C\) 给定后只
增加除子因子。

结合方程~\eqref{eq:inverse-first-separated}、加权 Cauchy 不等式
与旧标签纤维界，得到第一次
变换的显式输出：
\begin{equation}\label{eq:inverse-first-transform-output}
 \begin{split}
 |\mathcal T_{1,\sigma}|\ll{}&Z^{\pi_\beta+\pi_{\rm old}}
 \int_{\mathbb R^9}|\widehat{\mathfrak H}_{1,\sigma}(\boldsymbol t)|\\
 &\quad\times\prod_{i=1}^2
 \left(Z^{\kappa_i+9\eta/2}\mathcal S_{i,\sigma,\boldsymbol t}\right)^{1/2}
 \,d\boldsymbol t.
 \end{split}
\end{equation}
每个 \(\mathcal S_{i,\sigma,\boldsymbol t}\) 是方程~\eqref{eq:inverse-first-majorant}
中的正和：其行多项式带有
方程~\eqref{eq:inverse-first-polynomial} 中的
M\"obius 系数，由立方体导出的标签 \(J\) 被保留，而旧的 \(f,h\) 不再
出现在平均指标之中。以下变换应用于
这些正的行范数。

\paragraph{第二次 Poisson 变换。}
固定由第一个 Cauchy
不等式得到的两个正和之一。其多项式具有形式
\[
 P_i(y)=\sum_x c_i(x)\overline{\chi_x(y)}.
\]
系数 \(c_i(x)\) 由方程~\eqref{eq:first-poisson-column} 的其余因子组成，
其中 \(u_i=t'x\)、在 \(t'\) 处的刺点以及分离出的权重。它与
\(y\) 无关。特别地 \(x\) 无平方因子且
\[
 c_i(x)=0\quad\text{unless}\quad
 (x,CJ)=1,\quad (x,\operatorname{rad}q_0t')=1.
\]
刺点 \(\rho\) 以及 \(\chi_x(d_k)\) 的零点也
被保留。\(y\)-和使用了尺度为 \(Z^{H_{\rm use}}\) 的
固定形状光滑正球。

在其展开的平方中令
\[
 g'=(x_1,x_2),\qquad x_j=g'z_j,\qquad
 \widehat g=\log_Zq_{g'}.
\]
设 \(g\ge0\) 为其固定中心。
则 \(z_1,z_2\) 无平方因子且互素，且各自与
\(g'\) 互素。引理~\ref{lem:masked-poisson} 的行数据为
\begin{equation}\label{eq:second-masked-data}
 K=Z^{H_{\rm use}},\qquad
 \psi_2=\bar\chi_{z_1}\chi_{z_2},\qquad
 m_2=z_1z_2,\qquad R_{\rm mask}=g'.
\end{equation}
公因子 \(\bar\chi_{g'}(y)\chi_{g'}(y)\) 正是这个
掩码。该特征在 \(m_2=1\) 之外是本原的，因为其
两个非零局部幂具有不相交的支集。对 \(d_2\mid g'\)，
设 \(\theta\ge0\) 为其中心并记
\(\widehat\theta=\log_Zq_{d_2}\)。实际的导子长度为
\[
 \widehat L_2=\widehat x_1+\widehat x_2-2\widehat g,
 \qquad |\widehat L_2-2(s_i-g)|\le10\eta.
\]
最后一个不等式使用了每个副本的完整新 \(x\)-支集界
\(|\widehat x_a-s_i|\le4\eta\) 以及
\(|\widehat g-g|\le\eta\)。对 \(k_{\rm new}=d_kd_2k''\)，蕴含式
\(Z^{H_{\rm use}}q_{k''}/(q_{d_2}q_{m_2})\le Z^\tau\) 给出
\[
 \begin{split}
 \widehat{k_{\rm new}}
 &\le2(s_i-g)-H_{\rm use}+\delta+2\theta+13\eta+\tau\\
 &=M_c+\eta,
 \end{split}
\]
其中
\begin{equation}\label{eq:second-poisson-lengths}
 \begin{split}
 M_c&=2(s_i-g)-H_c+\delta+2\theta\\
    &=M-4\ell-2A_i-2t-2g+2\theta-2V+j,\qquad
 M_{\rm ch}=M_c+\eta.
 \end{split}
\end{equation}
\(13\eta\) 由来自导子的 \(10\eta\)、来自
\(d_k\) 的一个以及来自 \(d_2\) 的两个组成。减去
\(H_{\rm use}=H_c+12\eta+\tau\) 抵消了在这第二次 Poisson 比较中
所用的同一容差 \(\tau\)。因此不需要单独的行裁剪
或未记录的边界误差。

将原来的主贡献分离出来，用于下面的直接计数。
在任何非互素扩展或 Fourier 绝对化之前，仅在
真正的非主第二次 Poisson 表达式上保留外掩码
\[
 \mathcal B_2(k'')=1_{0<q_{d_kd_2k''}\le Z^{M_{\rm ch}}}.
\]
对每个被支集的配对，其余集包含在刚刚考虑的实际比值
尾部中，因此方程~\eqref{eq:inverse-raw-tail} 以
下面固定的阶将其丢弃。在掩码内部保留完整的光滑核。
掩码依赖于 \(d_k,d_2,k''\) 及固定扇区，而不依赖于剩余的
\(z_1,z_2\)。若 \(M_{\rm ch}<0\)，非零球为空，同样的
尾部比较丢弃每一个非零频率。在保留的非空
通道中 \(M_{\rm ch}\ge0\)。

实际的第二个前因子具有精确的侧分解
\[
 Z^{H_{\rm use}-\widehat\theta-\widehat L_2/2}
       =Z^{\widehat\lambda_1/2}Z^{\widehat\lambda_2/2},
 \qquad \widehat\lambda_a=H_{\rm use}-\widehat\theta-\widehat x_a+
                           \widehat g.
\]
对 \(\lambda_c=H_c-\theta-(s_i-g)\)，同样的支集界给出
\begin{equation}\label{eq:inverse-second-prefactor-error}
 \widehat\lambda_a\le\lambda_c+18\eta+\tau.
\end{equation}
这里 \(18=12+1+4+1\)。如同第一次分解，所有正规化的实
逆根将保留在联合光滑剖面中直至 Fourier
反演；它们也不会被追加到最终的子测试中。
所有原系数掩码都保留。特别地
\[
 (g',CJ)=1,\qquad (z_j,g'CJ)=1
\]
在其非零支集上成立。

主情形为 \(z_1=z_2=1\)，等价地 \(x_1=x_2\)。
其计数的名义恒等式为
\begin{equation}\label{eq:second-principal-exponent}
 \kappa_i+H_c+s_i+B+t+\ell+R/2=F-B-j,
 \qquad F=N+V=r+3\ell+V.
\end{equation}
实际上界计数加上来自第一个前因子的 \(9\eta/2\)、
来自行尺度 \(H_{\rm use}\) 的 \(12\eta+\tau\)、来自
对角 \(x\)-窗口的 \(4\eta\)、来自 \(C,t'\) 的各 \(\eta\)，以及来自
合并 \(q_0,J\) 计数的 \(7\eta/2\)。它们的和为 \(26\eta+\tau\)。以
聚合局部损失 \(\pi\) 计，该主贡献因此
至多为 \(Z^{F-B-j+26\eta+\tau+\pi}\le Z^{F+26\eta+\tau+\pi}\)。
任何被恢复到形式表达式中的非零主频率都带有
相同的 \(\mathcal B_2\) 掩码；其绝对和由引理~\ref{lem:masked-poisson}
中的完整主恢复所界定，因为
\(Z^{H_{\rm use}}\ge1\)。

对于非主项，我们现在识别新的系数
类。对无平方因子的 \(z\)，本原
Gauss 和的复共轭与方程~\eqref{eq:signal} 给出
\[
 \gamma_{-1}(z)=\chi_z(-1)\overline{\gamma_1(z)},\qquad
 \mu(z)\gamma_{-1}(z)
 =\bar\alpha(z)\gamma_2(z)\chi_z(-1)\overline{G(z)}.
\]
原始的第二个因子为
\(\mu(z)\gamma_1(z)=\overline{\bar\alpha(z)\gamma_2(z)}G(z)\)；
写在共轭的第二个
多项式内的对应因子具有射线因子 \(\overline{G(z)}\)。两个本原
根的 CRT 交叉因子为 \(\mathcal R(z_1,z_2)\)。因此它们合并的射线因子为
\[
 \mathfrak G(a,b)=\chi_a(-1)\overline{G(a)}G(b)\mathcal R(a,b)
       =G(ba^{-1}).
\]
最后一个等式来自
方程~\eqref{eq:G-quadratic-refinement} 以及
\(\mathcal R(a,a)=\chi_a(-1)\)，先在素数上由
方程~\eqref{eq:fixed-gauss-phases} 得到，然后乘性地推广。
此式中的所有商都在固定有限射线群中。
特别地，
\[
 \mathfrak G(v'n_1,v'n_2)=\mathfrak G(n_1,n_2).
\]
公乘性特征 \(\nu_i(v')\) 也在
两边之间抵消。

从对所有 \(d_k,d_2\) 公共的一个划分插入 \(q_{k''}\) 的二分块。
掩码 \(\mathcal B_2\) 蕴含 \(q_{k''}\le Z^{M_{\rm ch}}\)，因此在
有界保留区间内其
数目是对数的。设 \(h_2\) 为
这样一个裸频率二分块的中心。我们现在记录要分离的精确
分量。固定第一个模式与选定的侧 \(i\)，
记
\[
 H_o=CJ,\qquad r_0=\operatorname{rad}q_0t',\qquad d_0=d_k,
 \qquad a_0(z)=\bar\alpha(z)\gamma_2(z),\qquad
 \mathcal K_2=\mathcal F\Phi_+,
\]
并保留来自
方程~\eqref{eq:second-masked-data} 的 \(K=Z^{H_{\rm use}}\)。对无平方因子的自变量令
\begin{equation}\label{eq:inverse-second-outer-coefficients}
 \begin{split}
 B_o(g')&=\mu(g')\nu_i(g')\rho(g')1_{(g',r_0)=1}
              \chi_{g'}(H_o)^4\chi_{g'}(d_0),\\
 L_{o,d_2,k''}(z)&=a_0(z)\nu_i(z)\rho(z)1_{(z,r_0)=1}
       \chi_z(H_o)^4\chi_z(d_0)\overline{\chi_z(d_2)}\chi_z(k'').
 \end{split}
\end{equation}
\(P_i\) 中旧的分离测试记为 \(W_i\)，因此
\(W_i(x)=\omega_{x,i}(x)x^{\varepsilon_iit_{x_i}}\)。
设 \(\psi_g,\psi_{d_2},\psi_{k''}\) 表示在固定 \(g',d_2,k''\) 二分块上
保留的各个截断，取其正规化
范数处的值。对于这个初步模式 \(\sigma_0\)，方程~\eqref{eq:inverse-first-majorant}
中带掩码的
主恢复非零分量恰好为
\begin{equation}\label{eq:inverse-second-formal-component}
 \begin{split}
 \mathcal T_{2,\sigma_0,i}={}&
 \sum_{\substack{o\in\mathfrak O_\sigma,\ g'\ {\rm sf}\\
d_2\mid g',\ k''\ne0}}
 w_o^+\psi_g\psi_{d_2}\psi_{k''}\mu(d_2)\mathcal B_2(k'')|B_o(g')|^2\\
 &\times\sum_{\substack{z_1,z_2\ {\rm sf}\\
(z_1z_2,g')=1}}
 1_{(z_1,z_2)=1}L_{o,d_2,k''}(z_1)\overline{L_{o,d_2,k''}(z_2)}
       G([z_2][z_1]^{-1})\\
 &\times\mathfrak d_{I_i}(g'z_1)\overline{\mathfrak d_{I_i}(g'z_2)}
       W_i(q_{g'}q_{z_1}/Z^{s_i})
       \overline{W_i(q_{g'}q_{z_2}/Z^{s_i})}\\
 &\times\frac{K}{q_{d_2}\sqrt{q_{z_1}q_{z_2}}}
       \mathcal K_2\!\left(\frac{Kq_{k''}}{q_{d_2}q_{z_1}q_{z_2}}\right).
 \end{split}
\end{equation}
这里括号表示固定有限射线群中的类。的确，\(c_i(g'z)\) 的
算术部分分解为 \(B_o(g')\) 及其
在 \((g',z)=1\) 上的剩余 \(\mu(z)\) 系数。公行因子
是在 \(g'\) 处的掩码，其 Poisson 展开给出一个除子
\(d_2\)，而上面的信号计算给出所示的两个
\(L\) 因子与射线商。被恢复的主配对
\(z_1=z_2=1\) 由 \(G(1)=1\) 相合。其单独的代价，以及在此公式之前
移除的原始补集，已经得到界定。
因此这是对正
优函数指定分量的等式，而不是对原始带符号块的等式。

用其所示的无商 \(L\) 因子、固定射线商、形式乘积范数、
完整核以及不变的 \(\mathcal B_2\)，在
所有各自容许的无平方因子 \(z_1,z_2\) 上定义方程~\eqref{eq:inverse-second-formal-component}
的被加项。它在互素配对上与
真正的公式一致。在任何逐因子估计之前插入完整恒等式
\[
 1_{(z_1,z_2)=1}=\sum_{v'\mid(z_1,z_2)}\mu(v'),\qquad z_a=v'n_a,
\]
设 \(v\ge0\) 为
无平方因子 \(v'\)-二分块的中心，\(\widehat v=\log_Zq_{v'}\)。
\(n_a\) 无平方因子，且 \((v',n_a)=1\)；它们不必互相
互素。在此扩展上不使用本原 Poisson 公式。
对无平方因子的 \(v'\) 令
\[
 D_o(v';d_2,k'')=a_0(v')\nu_i(v')\rho(v')1_{(v',r_0)=1}
   \chi_{v'}(H_o)^4\chi_{v'}(d_0)\overline{\chi_{v'}(d_2)}\chi_{v'}(k'').
\]
在对无平方因子且互素的 \(v',n\) 上，精确提取为
\begin{equation}\label{eq:inverse-second-extraction}
 \begin{split}
 L_{o,d_2,k''}(v'n)={}&D_o(v';d_2,k'')a_0(n)\nu_i(n)\rho(n)1_{(n,r_0)=1}\\
 &\times\chi_n(H_o)^4\chi_n(d_0)\overline{\chi_n(d_2)}
          \chi_n(k'')\chi_n(v')^4.
 \end{split}
\end{equation}
这使用了 \(a_0(v'n)=a_0(v')a_0(n)\chi_n(v')^4\)。在重叠处，
右端通过 \(\chi_n(v')^4\) 定义该扩展为零，
而不去计算非无平方因子的 Gauss 和。两个公因子
给出 \(|D_o(v';d_2,k'')|^2\)，其中包括零点 \((v',k'')=1\) 及其
固定刺点、\(d_0\) 以及标签零点。它们通过 Cauchy 保留在外
权重中。旧的最大公因子还留下 \(1_{(v',g')=1}\)。
由于 \(D_o\) 在 \((v',CJ)>1\) 上为零，我们等价地保留
显式外因子 \(1_{(v',g'CJ)=1}\)，它重复了那个标签零点。

令 \(r_g=g'/d_2\)、\(k_{\rm new}=d_kd_2k''\) 以及
\(f_{\rm new}=JCd_2v'\)。所有关于 \(n\) 的变动因子满足
完整的零扩展恒等式
\begin{equation}\label{eq:inverse-moving-characters}
 \chi_n(d_k)\chi_n(k'')\overline{\chi_n(d_2)}
       \chi_n(CJ)^4\chi_n(v')^4
 =\chi_n(k_{\rm new})\chi_n(f_{\rm new})^4.
\end{equation}
它使用了 \(\overline{\chi_n(d_2)}=\chi_n(d_2)\chi_n(d_2)^4\)，在
非单位元上也成立。与 \(g'=d_2r_g\) 的剩余最大公因子分解为在 \(d_2\) 处重复的
零点以及在 \(r_g\) 处的固定刺点。在 \(q_0,t'\) 处的旧刺点
保持不变。若 \(\mathcal C_{\rm ray}\) 表示 \(G\) 通过其分解的固定
群，令
\[
 \widehat G(\vartheta)=|\mathcal C_{\rm ray}|^{-1}
       \sum_{c\in\mathcal C_{\rm ray}}G(c)\overline{\vartheta(c)}.
\]
则
\[
 G([n_2][n_1]^{-1})=\sum_{\vartheta\in\widehat{\mathcal C}_{\rm ray}}
   \widehat G(\vartheta)\overline{\vartheta(n_1)}\vartheta(n_2).
\]
因此在固定射线被加项中的两个新多项式使用相同的
\(\nu_i\overline\vartheta\)，而其唯一有界系数
\(\widehat G(\vartheta)\) 保持在外层。

\paragraph{新的典范数据及其容许性。}
第二次变换已将变动特征恢复到典范
形式。我们现在识别哪些数据将被固定、哪些将保持
平均。这一区分也决定了子的刺点。
对两侧令 \(\nu_a'=\nu_i\overline\vartheta\) 并定义
\begin{equation}\label{eq:canonical-child-data}
 \begin{gathered}
 \gamma=(q_0,t',r_g),\qquad r_g=g'/d_2,\\
 k_{\rm new}=d_kd_2k'',\qquad f_{\rm new}=JCd_2v',\\
 \rho'(n)=\rho_\gamma(n)=\rho(n)1_{(n,\operatorname{rad}q_0t'r_g)=1}.
 \end{gathered}
\end{equation}
在非零系数支集上，\(J,C,d_2,v'\) 无平方因子
且两两互素：\((C,J)=1\) 属于 \(\mathfrak O_\sigma\)，
\(|B_o(g')|^2\) 迫使 \((g',CJ)=1\)、\(d_2\mid g'\)，而 \(D_o\) 被
保留的零点连同旧的最大公因子限制给出
\((v',g'CJ)=1\)。因此 \(f_{\rm new}\) 在该支集上无平方因子。
这些公因子通过 Cauchy 保留在外权重中。

三元组 \(\gamma\) 将在子行和标签求和之前被固定。
其刺点 \(\rho_\gamma\) 因此独立于这两个
平均变量。理想 \(J,C,d_2,v'\) 保留在新的
平均标签中；它们不被计作额外的固定标签。

在每个选定侧的
副本上，以有序列表 \((d_2,r_g,v';n_a)\) 应用方程~\eqref{eq:inverse-slot-priority}。
设 \(I_a'\) 为副本 \(a\) 中在 \(n_a\) 上保留的选定侧槽位的
子集。被指定的素数及其优先级掩码保持在外层。
整除 \(n\) 的存活素数
已由四次幂零点与 \(JCd_2v'\) 互素，并
由固定刺点与 \(\operatorname{rad}q_0,t',r_g\) 互素。这些也
提供了所有更早的存活排除，因为
\(\operatorname{supp}\mathfrak B\subset\operatorname{supp}(Jq_0)\)。
存活的标记因此是一个原有的乘积形式子集。
一个被描述为指派给 \(J\) 的旧槽位，只是由这个导出的支集对
旧 \(b_a,\mathfrak A_a\) 指派的重分组，而不是
一个新的独立素数选择。

将 \(\mathcal B_2\) 保持在完整核
与旧窗口的 Fourier 分离之外。设 \(I_g,I_d,I_v,I_k\) 为
当前各个二分块截断 \(\psi_g,\psi_{d_2},\psi_v,\psi_{k''}\) 的完整支集。
选取新的截断 \(\omega_{n,a}\)，使其在
\(W_i\) 的支集除以 \(I_gI_v\) 上等于一，并在区间
\([a,b]\) 内具有固定的完整支集，其中 \(b\ge1\)。选取一个非负的 \(\omega_f\)，使其在
\(I_JI_CI_dI_v\) 上等于一，并具有固定的完整支集 \(I_f\)。将这些
完整支集，以及在其上重构的形式乘积 \(g'v'n\)，
纳入上述指定的有限窗口族 \(\mathcal I_D\)。
这些完整的全新子窗口给出形式中心与支集界
\begin{equation}\label{eq:canonical-child-lengths}
 \begin{gathered}
 N_c=r-A_i-B-t-g-v,\qquad V_c=B+\theta+v+j,\qquad F_c=N_c+V_c,\\
 |\widehat n_{\rm child}-N_c|\le6\eta,\qquad
 |\widehat f_{\rm new}-V_c|\le4\eta.
 \end{gathered}
\end{equation}
在因式分解表达式上，这些界来自精确乘积
\(n_i=\mathfrak A_iCt'g'v'n_{\rm child}\) 与
\(f_{\rm new}=JCd_2v'\)。完整的全新子截断与完整扩大的
标签窗口包含在 \(\mathcal I_D\) 中，因此当稍后由
正性加入独立项时，同样的界直接在其支集上成立。那些加入的项
并不被断言来自父因式分解。在当前分离过程中，在完整全新截断上
重构的形式乘积 \(x=g'v'n\) 也包含在
\(\mathcal I_D\) 中；因此第二次前因子中所用的 \(4\eta\) 界
在整个分离侧上成立。

保持非负标签中心 \(V_c\)，不向上取整。
仅当保留的行分量非零时才进入子；此时
外球有 \(M_{\rm ch}\ge0\)。这个门控独立于任何
结构外指标集，后者即使在
\(\mathcal B_2\) 恒为零时也可能包含零权元组。类似地，若一个完整的全新
子列窗口不包含无平方因子理想，则其多项式为零
且该分量被省略。若 \(N_c<0\) 且存在子列，则其范数至少为一，且
方程~\eqref{eq:canonical-child-lengths} 给出 \(-6\eta\le N_c<0\)。
对一个被保留的非空子定义
\begin{equation}\label{eq:inverse-child-clipping}
 N_{\rm ch}=\max(0,N_c),\qquad \delta_N=N_{\rm ch}-N_c\in[0,6\eta],
 \qquad F_{\rm ch}=N_{\rm ch}+V_c=F_c+\delta_N.
\end{equation}
这是测试的有界改变，而不是 \(Z^\eta\) 宽度的支集
扩大。使用上面选取的固定包围支集区间 \([a,b]\)。若
\(N_c<0\)，非空性给出 \(Z^{\delta_N}\le b\)；若 \(N_c\ge0\)，
则 \(\delta_N=0\)，同样的界由 \(b\ge1\) 得到。
因此在两种情形下都有 \(1\le Z^{\delta_N}\le b\)。新的环形截断
\(\omega_{n,a}(Z^{\delta_N}y)\) 在 \([a/b,b]\) 中有支集且具有一致
有界的 Euler 半范数。它们的尺度只依赖于固定的扇区
中心，而不依赖于当前的行或标签。这些裁剪
支集的完整并集包含在 \(\mathcal I_D\) 中。

通过 \(v'\)-二分块、
第二个射线被加项与槽位分支固定 \(\sigma_0\) 的一个细化 \(\sigma\)。
令 \(\Xi_\sigma\) 由
\[
 \xi=(o,g'\ {\rm sf},d_2\mid g',v'\ {\rm sf},k''\ne0;
               \text{新指派的槽位素数})
\]
组成，带有保留的各个支集与槽位分支。它既不含
旧的 \(f,h\)，也不含当前的列指标 \(n_1,n_2\)。
定义 \(R_1,E,J,q_0\) 的恒等式继承自 \(o\)，
而 \(\gamma,k_{\rm new},f_{\rm new}\) 是 \(\xi\) 的所示函数。
该集合可以在施加 \(\mathcal B_2\) 与
\((v',g'CJ)=1\) 之前取定，将这些掩码保留在带符号的外权重中。令
\(\mathcal A_{2,\sigma}(\xi)\) 为所有新指派的
原槽位系数及其优先级掩码的乘积，在
侧二上取共轭。

设 \(\mathfrak X_\sigma^{\rm src}\) 为
这个变换分量的结构外集合，在加入任何独立子行
或标签之前。它保留
\(o\in\mathfrak O_\sigma\)，特别是精确重构
\(b_1b_2=q_0^2J_2^2s\)、导出的 \(J=sJ_2,q_0\)，以及完整的
固定 \(b_1,b_2,t'\) 二分块支集；它还保留固定的
\(g',d_2\) 支集、\(d_2\mid g'\)，以及它的其他外算术
与槽位条件。零权元组可以被保留，而单位
Fourier 相位在定义这个结构集时被忽略。用
下式定义不同三元组的集合，不含见证重数：
\begin{equation}\label{eq:inverse-source-triples}
 \Gamma_\sigma=
 \{(q_0,t',r_g=g'/d_2):\xi\in\mathfrak X_\sigma^{\rm src}\}.
\end{equation}
这个投影在所有可能的当前行与
标签上一次取定，在固定任何 \(k_{\rm new},f_{\rm new}\) 或 Fourier
模式之前。每个非零变换项都映入其中。每个元素
都有一个源见证，给出
\(\widehat c\le\ell+\eta\)、\(\widehat t\le t+\eta\) 以及
\(\log_Zq_{r_g}\le g-\theta+2\eta\)。这些界中的第一个
将为子刺点所需要；用于计数的包围范数球
并不蕴含它。上面陈述的非空行与子列门控
保持与这个结构投影分离。

令
\[
 D_c=\ell+A_i+t+g-\theta+V.
\]
由于 \(r=N-3\ell\) 是联系到固定父的一个精确定义，
在方程~\eqref{eq:canonical-child-lengths} 与
\eqref{eq:second-poisson-lengths} 中代入给出精确的形式恒等式
\begin{equation}\label{eq:canonical-transition}
 \begin{gathered}
 F_c=F-D_c-2\ell+j,\qquad M_c=M-2D_c-2\ell+j,\\
 F_c-M_c=F-M+D_c,\\
 4F_c-3M_c=4F-3M+2(A_i+t+g-\theta+V)+j.
 \end{gathered}
\end{equation}
这些恒等式涉及与
方程~\eqref{eq:inverse-fixed-invariant} 相同的固定 \(F=N+V\)，而不是实际的父范数。

对每个 \(\gamma\in\Gamma_\sigma\)，任取一个源见证。
其实际的整除性 \(d_2\mid g'\) 与
方程~\eqref{eq:inverse-cube-actual} 给出
\[
 \widehat\theta\le\widehat g,\qquad
 \widehat j\le2\widehat\ell_{\rm pair},\qquad
 \widehat c\le\widehat\ell_{\rm pair}.
\]
因此这里所需的唯一实际到中心的不等式为
\begin{equation}\label{eq:inverse-center-inequalities}
 g-\theta\ge-2\eta,\qquad j\le2\ell+3\eta,\qquad
 \widehat c\le\ell+\eta.
\end{equation}
前两个并不以零误差成立。特别地，
第三个不等式对其源见证的每个固定三元组成立，
尽管它不必在整个包围计数
球内成立。方程~\eqref{eq:canonical-child-data} 中的新根式
具有实际对数范数 \(Q_{\rm new}=Q_\gamma\)，满足
\begin{equation}\label{eq:inverse-new-puncture}
 \begin{split}
 Q_{\rm new}&\le Q+\widehat c+\widehat t+\widehat g-\widehat\theta\\
 &\le Q+\ell+t+g-\theta+4\eta\le Q+D_c+4\eta.
 \end{split}
\end{equation}
这在其因子与旧刺点共有素数时也成立，
因为取根式只会减小范数。最后一个不等式
使用了 \(A_i,V\ge0\)。

在方程~\eqref{eq:canonical-transition} 中使用 \(F_{\rm ch}=F_c+\delta_N\)、\(M_{\rm ch}=M_c+\eta\) 以及
\(2(A_i+t+g-\theta+V)+j\ge-4\eta\)。两个子余量满足
\begin{equation}\label{eq:inverse-child-margins}
 \begin{aligned}
 F_{\rm ch}-M_{\rm ch}-Q_{\rm new}-z_0
   &\ge c_{\rm node}+\delta_N-5\eta\ge c_{\rm node}-5\eta,\\
 4F_{\rm ch}-3M_{\rm ch}-6z_0
   &\ge c_{\rm node}+4\delta_N-7\eta\ge c_{\rm node}-7\eta.
 \end{aligned}
\end{equation}
名义上界 \(z_0\) 不变，因为存活的标记是一个
子集。此外 \(D_c\ge\ell+V-2\eta\) 且
\(2\ell-j\ge-3\eta\)，于是
\begin{equation}\label{eq:inverse-row-decrease}
 M-M_{\rm ch}=2D_c+2\ell-j-\eta
       \ge2(\ell+V)-8\eta\ge2d-8\eta.
\end{equation}
对 \(\eta\le d/16\)，这至少为 \(3d/2\)，因此至少为 \(d\)。
同样的恒等式给出 \(F_c\le F-\ell-V+5\eta\)，故
\(F_{\rm ch}\le F+11\eta\)。由于其两个被加项都非负，
这在有限深度内界定了子 \(N_{\rm ch},V_c\)。

算术构造因此产生了具有更小行长度的容许子
区间。我们接下来用它们的
典范多项式来表示变换后的和。所需的 Fourier 密度对所有
三元组 \(\gamma\)、行和标签都是公共的；只有在那个恒等式之后我们才
应用 Cauchy 并控制重构重数。

\paragraph{第二次变换和的公共密度。}
在对任何实际外理想求和之前，使用六个独立坐标
\[
 (y_g,y_d,y_v,y_k,y_1,y_2)
 =\left(\frac{q_{g'}}{Z^g},\frac{q_{d_2}}{Z^\theta},
 \frac{q_{v'}}{Z^v},\frac{q_{k''}}{Z^{h_2}},
 \frac{q_{n_1}}{Z^{N_c}},\frac{q_{n_2}}{Z^{N_c}}\right).
\]
使用已经选取的新截断 \(\omega_{n,a},\omega_f\)。
将四个当前二分块截断
与 \(\omega_f\) 保留在反演之外，而将 \(\omega_{n,a}\) 保留在两个
列中。设 \(\Omega_{2,\alpha}\) 为更大的截断，在相应四个二分块与两个列
截断的完整支集上等于一。这些完整支集、\(I_f\)，以及在其上的形式乘积
\(g'v'n\)，就是已经包含在 \(\mathcal I_D\) 中的窗口。

精确的第二根与核自变量为
\begin{equation}\label{eq:inverse-second-root-kernel}
 \begin{split}
 \frac{K}{q_{d_2}q_{v'}\sqrt{q_{n_1}q_{n_2}}}
 &=Z^{\lambda_c+12\eta+\tau}y_d^{-1}y_v^{-1}(y_1y_2)^{-1/2},\\
 \frac{Kq_{k''}}{q_{d_2}q_{v'}^2q_{n_1}q_{n_2}}
 &=A_{{\rm ker},2}\frac{y_k}{y_dy_v^2y_1y_2},\\
 A_{{\rm ker},2}&=Z^{H_{\rm use}+h_2-\theta-2v-2N_c}.
 \end{split}
\end{equation}
的确 \(N_c=s_i-g-v\) 且
\(\lambda_c=H_c-\theta-(s_i-g)\)。定义单一的环境剖面
\begin{equation}\label{eq:inverse-second-profile}
 \begin{split}
 \mathfrak H_{2,\sigma,\boldsymbol t}(\boldsymbol y)={}&Z^{-6\eta}
 \prod_\alpha\Omega_{2,\alpha}(y_\alpha)
 y_d^{-1}y_v^{-1}(y_1y_2)^{-1/2}\\
 &\times W_i(y_gy_vy_1)\overline{W_i(y_gy_vy_2)}
 \mathcal K_2\!\left(A_{{\rm ker},2}\frac{y_k}{y_dy_v^2y_1y_2}\right).
 \end{split}
\end{equation}
其外的标量恰好为 \(Z^{\lambda_c+18\eta+\tau}\)。
完整的 \(q_{d_2}\) 分母与两个正规化的实列根
都在这个剖面中，且没有这样的根被追加到最终测试中。

对 \(\boldsymbol\zeta=(\zeta_g,\zeta_d,\zeta_v,\zeta_k,
\zeta_1,\zeta_2)\)，令 \(z_n=q_n/Z^{N_{\rm ch}}\) 并定义
\[
 W_a'(z)=\omega_{n,a}(Z^{\delta_N}z)z^{\varepsilon_a i\zeta_a},
 \qquad \nu_a'=\nu_i\overline\vartheta,\qquad
 \mathfrak d_a'=\mathfrak d_{I_a'}.
\]
则任一新 Cauchy 侧上的多项式恰好为
\begin{equation}\label{eq:canonical-child-polynomial}
 \begin{split}
 Q_{a,\gamma,\boldsymbol\zeta}(k_{\rm new},f_{\rm new})={}&
 \sum_{n\ {\rm sf}}\bar\alpha(n)\gamma_2(n)\nu_a'(n)\rho'(n)
 \chi_n(k_{\rm new})\\
 &\times\chi_n(f_{\rm new})^4
 \mathfrak d_a'(n)W_a'(q_n/Z^{N_{\rm ch}}).
 \end{split}
\end{equation}
在 \(q_0,t',g'/d_2\) 与有限射线扇区固定之后，
\(\nu_a'\) 与 \(\rho'\) 独立于 \(J,C,d_2,v'\) 以及
新行。关于 \(n\) 的唯一新特征是显式
所示的行因子与四次幂因子；所有涉及
旧 \(b_i,\mathfrak A_i\) 的其他相位都被放入
方程~\eqref{eq:retained-cube-label} 或旧行
\(\widetilde h\) 中。分离出的环形权重 \(W_a'\) 也
独立于新行与标签。
完整的剩余外权重为
\begin{equation}\label{eq:inverse-second-outer-weight}
 \begin{split}
 V_\sigma(\xi;\boldsymbol\zeta)={}&w_o^+\mu(d_2)\mu(v')\widehat G(\vartheta)
 |B_o(g')|^2|D_o(v';d_2,k'')|^2
 1_{(v',g'CJ)=1}\mathcal B_2(k'')\\
 &\times\mathcal A_{2,\sigma}(\xi)
 \psi_g(y_g)\psi_{d_2}(y_d)\psi_v(y_v)\psi_{k''}(y_k)
 \omega_f(q_{f_{\rm new}}/Z^{V_c})\\
 &\times y_g^{i\zeta_g}y_d^{i\zeta_d}y_v^{i\zeta_v}y_k^{i\zeta_k}
       Z^{i\delta_N(\zeta_1+\zeta_2)}.
 \end{split}
\end{equation}
所有旧的纯算术与优先级掩码都保留在
\(\mathfrak O_\sigma\) 中，带有其非负的旧系数权重
\(w_o^+\)。所示的两个公共模保留其他每一个
公共系数零点，包括依赖于 \(k''\) 的那个。
有一个 \(\mu(d_2)\)、一个 \(\mu(v')\) 以及一个射线系数。
标签截断作为等式插入到原乘积
支集上并保持在外层。
